
\section{\BookRecordIndex{T0166}横截性}\label{sec:3-2}
\zhindex{hengjiexing}{横截性}

定理~\ref{thm:3-1-6} 和定理~\ref{thm:3-1-8} 的证明以 Smale
\srcite{381} 提出的无穷维\BookRecordIndex{T0296}萨德定理（见定理 A.5.1）为基础。粗略地说，设 $X,Y$
是\BookRecordIndex{T0213}可分\BookRecordIndex{T0053}巴拿赫流形，并且
\[
  F:X\longrightarrow Y
\]
是指标为 $k$ 的弗雷德霍姆映射；只要 $F$ 具有足够高的可微性（至少应为
$C^{k+1}$），$F$ 的\BookRecordIndex{T0409}正则值集 $Y_{\mathrm{reg}}$ 就是\BookRecordIndex{T0300}剩余集。与有限维情形一样，若对
每个满足 $F(x)=y$ 的 $x\in X$，\BookRecordIndex{T0346}线性化算子
\[
  dF(x):T_xX\longrightarrow T_yY
\]
都是满射，就称 $y\in Y$ 是\BookRecordIndex{T0409}正则值。此时由\BookRecordIndex{T0380}隐函数定理可知，对每个
$y\in Y_{\mathrm{reg}}$，$F^{-1}(y)$ 都是 $k$ 维流形（定理 A.3.3）。

为了在光滑范畴中得到结果，我们采取如下策略：对每个 $\ell$，先证明 $C^\ell$ 类正则
\BookRecordIndex{T0088}殆复结构的集合在 $C^\ell$ 拓扑下是通有的（即剩余的），再对所有 $\ell$ 取这些集合的
交。这一方法源于 Taubes。下文用 $\mathcal J^\ell$ 同时表示所有 $C^\ell$ 类
$\omega$-\BookRecordIndex{T0348}相容\BookRecordIndex{T0088}殆复结构组成的空间 $\mathcal J^\ell(M,\omega)$，或所有 $C^\ell$ 类
$\omega$-\BookRecordIndex{T0372}驯服\BookRecordIndex{T0088}殆复结构组成的空间 $\mathcal J_\tau^\ell(M,\omega)$。下面的
命题~\ref{prop:3-2-1} 对这两种情形都成立。

\paragraph{\BookRecordIndex{T0268}普适模空间。}
\BookRecordIndex{T0268}普适模空间
\[
  \mathcal M^*(A,\Sigma;\mathcal J^\ell)
  :=\bigl\{(u,J)\bigm|J\in\mathcal J^\ell,
  u\in\mathcal M^*(A,\Sigma;J)\bigr\}
\]
由所有表示同调类 $A$ 的单纯 \BookRecordIndex{T0003}$J$-全纯曲线 $u:\Sigma\to M$ 组成，其中 $J$ 在
$\mathcal J^\ell$ 中变化。回忆命题~\ref{prop:3-1-10}：当 $J$ 属于 $C^\ell$ 时，每条
\BookRecordIndex{T0003}$J$-全纯曲线都属于 $C^\ell$。因此，对任意 $p>2$ 及任意
$k\in\{1,\ldots,\ell\}$，可把 $\mathcal M^*(A,\Sigma;\mathcal J^\ell)$ 看成
$\mathcal B^{k,p}\times\mathcal J^\ell$ 的子集；这里与上文一样，
$\mathcal B^{k,p}$ 表示所有代表 $A$ 的 $W^{k,p}$ 映射 $u:\Sigma\to M$ 组成的空间。
在 $k=1$ 时必须小心，因为 $W^{1,p}$ 函数未必连续可微。如果连续映射
$u:\Sigma\to M$ 存在一点 $z\in\Sigma$ 和常数 $\delta>0$，使得对每个
$\zeta\in\Sigma$ 都有
\[
  d_M\bigl(u(z),u(\zeta)\bigr)
  \ge \delta d_\Sigma(z,\zeta),
\]
就称 $u$\emph{\BookRecordIndex{T0397}在某处单射}。这一条件与写出它时在 $M$ 和 $\Sigma$ 上选取的度量无关。
还要注意，当 $u$ 连续可微时，这一条件与第~\ref{sec:2-5} 节引入的\BookRecordIndex{T0397}在某处单射概念一致。

参数空间 $\mathcal J^\ell$ 是光滑\BookRecordIndex{T0213}可分\BookRecordIndex{T0053}巴拿赫流形。在 $\omega$-\BookRecordIndex{T0348}相容情形，
$J$ 处的切空间 $T_J\mathcal J^\ell$ 由丛 $\End(TM,J,\omega)$ 的 $C^\ell$ 截面
$Y$ 组成；该丛在 $x\in M$ 处的纤维是所有满足
\[
  YJ+JY=0,
  \qquad
  \omega(Yv,w)+\omega(v,Yw)=0
\]
的线性映射 $Y:T_xM\to T_xM$ 所成的空间。第一式是恒等式 $J^2=-\mathbf 1$ 的微分，
第二式来自\BookRecordIndex{T0348}相容性条件。等价地，也可把这些条件写成
\[
  Y=Y^*=JYJ,
\]
其中 $Y^*$ 表示关于度量 $g_J$ 的伴随算子。丛
$\End(TM,J,\omega)\to M$ 的 $C^\ell$ 截面空间是巴拿赫空间，并通过
\[
  Y\longmapsto J\exp(-JY)
\]
给出 $\mathcal J^\ell$ 的局部模型。相应的 $C^\infty$ 截面空间不是巴拿赫空间，
而只是弗雷歇空间。我们的约定是：没有上标的空间均由 $C^\infty$ 光滑元素组成。

\begin{proposition}\label{prop:3-2-1}
固定同调类 $A\in H_2(M;\mathbb Z)$、整数 $\ell\ge2$、实数 $p>2$ 以及整数
$k\in\{1,\ldots,\ell\}$。则\BookRecordIndex{T0268}普适模空间
$\mathcal M^*(A,\Sigma;\mathcal J^\ell)$ 是
$\mathcal B^{k,p}\times\mathcal J^\ell$ 的一个\BookRecordIndex{T0213}可分 $C^{\ell-k}$ 巴拿赫子流形。
\end{proposition}

下面将看到，这一命题成立是因为切空间 $T_J\mathcal J^\ell$ 足够大。特别地，下一个
引理表明，它在每一点 $x\in M$ 都传递地作用于 $T_xM$。

\begin{lemma}\label{lem:3-2-2}
设 $\xi,\eta\in\mathbb R^{2n}$ 是两个非零向量。则存在矩阵
$Y\in\mathbb R^{2n\times2n}$，使得
\[
  Y=Y^T=J_0YJ_0,
  \qquad
  Y\xi=\eta.
\]
\end{lemma}

\begin{proof}
$Y$ 的一个显式公式为
\begin{align*}
  Y:={}&\frac1{|\xi|^2}
  \Bigl(\eta\xi^T+\xi\eta^T
  +J_0(\eta\xi^T+\xi\eta^T)J_0\Bigr)\\
  &-\frac1{|\xi|^4}\Bigl(
  \langle\eta,\xi\rangle(\xi\xi^T+J_0\xi\xi^TJ_0)
  +\langle\eta,J_0\xi\rangle(J_0\xi\xi^T-\xi\xi^TJ_0)
  \Bigr).
\end{align*}
这个公式是分步构造的。第一项给出一个把 $\xi$ 映到 $\eta$ 的矩阵；第二项使它对称，
第三、第四项则使等式 $Y=J_0YJ_0$ 成立。然而这样得到的矩阵不再把 $\xi$ 映到
$\eta$，最后四项正好对此作出校正。
\end{proof}

\begin{proof}[命题~\ref{prop:3-2-1} 的证明]
考虑 $C^{\ell-k}$ 巴拿赫向量丛
\[
  \mathcal E^{k-1,p}\longrightarrow
  \mathcal B^{k,p}\times\mathcal J^\ell,
\]
它在 $(u,J)$ 上的纤维为
\[
  \mathcal E_{(u,J)}^{k-1,p}
  :=W^{k-1,p}\bigl(\Sigma,
  \Lambda^{0,1}\otimes_Ju^*TM\bigr),
\]
即 $\Sigma$ 上取值于 $u^*TM$ 的 $W^{k-1,p}$ 类 $J$-反线性 $1$-形式所成的空间。
可以如下构造 $\mathcal E^{k-1,p}$ 的 $C^{\ell-k}$ \BookRecordIndex{T0053}巴拿赫流形结构。固定 $M$ 上任意
光滑度量，通过 $\xi\mapsto\exp_u(\xi)$，把 $\mathcal B^{k,p}$ 中光滑元素
$u:\Sigma\to M$ 的一个邻域 $\mathcal N(u)$ 与巴拿赫空间
$W^{k,p}(\Sigma,u^*TM)$ 中零点的一个邻域等同起来。这样就得到
$\mathcal B^{k,p}$ 上具有光滑转移映射的坐标卡。再用\BookRecordIndex{T0052}埃尔米特联络
$\widetilde\nabla$，在这样的坐标卡上平凡化丛 $\mathcal E^{k-1,p}$。由于 $J$ 属于
$C^\ell$，联络 $\widetilde\nabla$ 及其\BookRecordIndex{T0263}平行移动映射属于 $C^{\ell-1}$。所以，即使把由
这些平凡化产生的转移映射微分 $\ell-k$ 次，所得映射在 $W^{k-1,p}$ 中仍有定义。这就在
切片 $\mathcal B^{k,p}\times\{J\}$ 的开集 $\mathcal N(u)$ 上构造了局部平凡化。

为了把它扩张到 $\mathcal N(u)\times\mathcal N(J)$ 这种形状的邻域，先利用同构
\[
  \Lambda^{0,1}\otimes_Ju^*TM
  \longrightarrow
  \Lambda^{0,1}\otimes_{J'}u^*TM,
  \qquad
  \alpha\longmapsto\frac12(\alpha+J'\circ\alpha\circ j),
\]
在一个适当的邻域 $\{u\}\times\mathcal N(J)$ 上平凡化，然后像刚才一样用\BookRecordIndex{T0263}平行移动，
在每个切片 $\mathcal N(u)\times\{J'\}$ 上扩张这一平凡化。这说明
$\mathcal E^{k-1,p}$ 是 $C^{\ell-k}$ 巴拿赫向量丛。

映射 $(u,J)\mapsto\bar\partial_J(u)$ 定义了丛
$\mathcal E^{k-1,p}\to\mathcal B^{k,p}\times\mathcal J^\ell$ 的一个
$C^{\ell-k}$ 截面。把这个截面记作
\[
  \mathcal F:\mathcal B^{k,p}\times\mathcal J^\ell
  \longrightarrow\mathcal E^{k-1,p},
  \qquad
  \mathcal F(u,J)=\bar\partial_J(u).
\]
我们须证明：只要 $u$ 单纯且 $\mathcal F(u,J)=0$，\BookRecordIndex{T0310}竖直微分
\begin{multline*}
  D\mathcal F(u,J):
  W^{k,p}(\Sigma,u^*TM)
  \times C^\ell\bigl(M,\End(TM,J,\omega)\bigr)\\
  \longrightarrow
  W^{k-1,p}\bigl(\Sigma,
  \Lambda^{0,1}\otimes_Ju^*TM\bigr)
\end{multline*}
就是满射。由命题~\ref{prop:3-1-1}，这一微分由公式
\[
  D\mathcal F(u,J)(\xi,Y)
  =D_u\xi+\frac12Y(u)\circ du\circ j
\]
给出。由于 $D_u$ 是\BookRecordIndex{T0128}弗雷德霍姆算子，$D\mathcal F(u,J)$ 的像是闭的；因此，当
$\bar\partial_J(u)=0$ 时，只须证明它的像稠密。

先考虑 $k=1$。若像不稠密，由 Hahn--Banach 定理，存在非零截面
\[
  \eta\in L^q\bigl(\Lambda^{0,1}\otimes_Ju^*TM\bigr),
  \qquad \frac1p+\frac1q=1,
\]
使它消去 $D\mathcal F(u,J)$ 的像。也就是说，对每个
$\xi\in W^{1,p}(\Sigma,u^*TM)$，都有
\begin{equation}\label{eq:3-2-1}
  \int_\Sigma\langle\eta,D_u\xi\rangle\,
  d\operatorname{vol}_\Sigma=0,
\end{equation}
并且对每个 $Y\in C^\ell(M,\End(TM,J,\omega))$，都有
\begin{equation}\label{eq:3-2-2}
  \int_\Sigma
  \langle\eta,Y(u)\circ du\circ j\rangle\,
  d\operatorname{vol}_\Sigma=0,
\end{equation}
其中 $j:=j_\Sigma$。由式~\eqref{eq:3-2-1} 和命题~\ref{prop:3-1-11} 可知，
$\eta$ 属于 $W^{1,p}$，并且 $D_u^*\eta=0$。

由于 $u$ 单纯，命题~\ref{prop:2-5-1} 断言 $u$ 的\BookRecordIndex{T0094}单射点集在 $\Sigma$ 中开且稠密。
取一个这样的\BookRecordIndex{T0094}单射点 $z_0\in\Sigma$，即
\[
  du(z_0)\ne0,
  \qquad
  u^{-1}\bigl(u(z_0)\bigr)=\{z_0\}.
\]
我们来证明 $\eta$ 在 $z_0$ 处为零。反设 $\eta(z_0)\ne0$。由于 $du(z_0)\ne0$，
由引理~\ref{lem:3-2-2}，存在
\[
  Y_0\in\End\bigl(T_{u(z_0)}M,J_{u(z_0)},\omega_{u(z_0)}\bigr),
\]
使得
\[
  \bigl\langle\eta(z_0),
  Y_0\circ du(z_0)\circ j(z_0)\bigr\rangle>0.
\]
任取满足 $Y(u(z_0))=Y_0$ 的截面
$Y\in C^\ell(M,\End(TM,J,\omega))$。则 $\Sigma$ 上的标量函数
\[
  \langle\eta,Y(u)\circ du\circ j\rangle
\]
在 $z_0$ 的某个开邻域 $V_0\subset\Sigma$ 中为正。由于 $z_0$ 是 $u$ 的\BookRecordIndex{T0094}单射点，紧集
$u(\Sigma\setminus V_0)$ 不含点 $u(z_0)$。因此存在 $u(z_0)$ 在 $M$ 中的开邻域
$U_0$，使得
\[
  u(\Sigma\setminus V_0)\cap U_0=\varnothing,
  \qquad\text{从而}\qquad
  u^{-1}(U_0)\subset V_0.
\]
现取支集含于 $U_0$ 且满足 $\beta(u(z_0))=1$ 的光滑\BookRecordIndex{T0188}截断函数
$\beta:M\to[0,1]$。于是函数
\[
  \beta(u)\langle\eta,Y(u)\circ du\circ j\rangle
\]
在 $\Sigma$ 上的支集含于 $V_0$，处处非负且在某处为正。因此，把
式~\eqref{eq:3-2-2} 中的 $Y$ 换成 $\beta Y$，其左端积分不为零，矛盾。这说明
$\eta(z_0)=0$。这一结论对 $u$ 的每个\BookRecordIndex{T0094}单射点 $z_0$ 都成立，而\BookRecordIndex{T0094}单射点集在 $\Sigma$ 中
稠密（命题~\ref{prop:2-5-1}），故 $\eta$ 几乎处处为零。又因 $\eta$ 连续，遂有
$\eta\equiv0$。所以 $D\mathcal F(u,J)$ 的像稠密，而在 $k=1$ 时它因而是满射。

再证明一般 $k$ 的满射性。给定
\[
  \eta\in W^{k-1,p}\bigl(\Sigma,
  \Lambda^{0,1}\otimes_Ju^*TM\bigr),
\]
由 $k=1$ 时的满射性，存在
\[
  (\xi,Y)\in W^{1,p}(\Sigma,u^*TM)
  \times C^\ell\bigl(M,\End(TM,J,\omega)\bigr)
\]
使 $D\mathcal F(u,J)(\xi,Y)=\eta$。上面的微分公式给出
\[
  D_u\xi=\eta-\frac12Y(u)\circ du\circ j\in W^{k-1,p}.
\]
由\BookRecordIndex{T0329}椭圆正则性，$\xi\in W^{k,p}(\Sigma,u^*TM)$（定理 C.2.3）。因此，对每个
$(u,J)\in\mathcal M^*(A,\Sigma;\mathcal J^\ell)$，$D\mathcal F(u,J)$ 都是满射。

因为 $D_u$ 是\BookRecordIndex{T0128}弗雷德霍姆算子，由引理 A.3.6 可知，$D\mathcal F(u,J)$ 有\BookRecordIndex{T0392}右逆。
无穷维\BookRecordIndex{T0380}隐函数定理（定理 A.3.3）因而说明，
$\mathcal M^*(A,\Sigma;\mathcal J^\ell)$ 是
$\mathcal B^{k,p}\times\mathcal J^\ell$ 的 $C^{\ell-k}$ 巴拿赫子流形。由于后者
\BookRecordIndex{T0213}可分，$\mathcal M^*(A,\Sigma;\mathcal J^\ell)$ 也\BookRecordIndex{T0213}可分。这就证明了
命题~\ref{prop:3-2-1}。
\end{proof}

\begin{remark}\label{rem:3-2-3}
命题~\ref{prop:3-2-1} 的断言对更一般的\BookRecordIndex{T0088}殆复结构集合 $\mathcal J$ 仍成立。例如，可以只
考虑所有在开集 $U\subset M$ 的补集上与给定\BookRecordIndex{T0088}殆复结构 $J_0$ 一致的 $\omega$-\BookRecordIndex{T0348}相容\BookRecordIndex{T0088}殆复
结构，只要对每个 $J$，每条代表 $A$ 的 \BookRecordIndex{T0003}$J$-全纯曲线 $u:\Sigma\to M$ 都经过 $U$。
不过在这种情形，证明需要定理~\ref{thm:2-3-4}。事实上，若 $D_u^*\eta=0$，则
\[
  0=D_uD_u^*\eta=\Delta\eta+\text{低阶项}.
\]
由 Aronszajn 定理，如果 $\eta$ 在某个开集上为零，就有 $\eta\equiv0$。现在，只要选取
支集含于 $U$ 的\BookRecordIndex{T0343}无穷小殆复结构 $Y$，就能保证命题~\ref{prop:3-2-1} 证明中的零化子
$\eta$ 在某个开集上为零。稍后在引理~\ref{lem:3-4-3} 和引理~\ref{lem:3-4-6} 中，
考察曲线在定义域中有限个
点上的取值受到约束的\BookRecordIndex{T0246}模空间时，我们还会用到这一想法。
\end{remark}

\paragraph{主要定理的证明。}
准备工作已经完成，现在可以证明本章的主要定理。

\begin{proof}[定理~\ref{thm:3-1-6}\textup{(i)} 的证明]
设 $J\in\mathcal J_{\mathrm{reg}}(A,\Sigma)$，并且
$u\in\mathcal M^*(A,\Sigma;J)$ 是 \BookRecordIndex{T0003}$J$-全纯曲线。由命题~\ref{prop:3-1-10}，$u$
光滑。给定整数 $k\ge1$ 和实数 $p>2$，考虑由式~\eqref{eq:3-1-3} 定义的映射
\[
  \mathcal F_u:W^{k,p}(\Sigma,u^*TM)
  \longrightarrow
  W^{k-1,p}\bigl(\Sigma,
  \Lambda^{0,1}\otimes_Ju^*TM\bigr).
\]
$\mathcal F_u^{-1}(0)$ 中零点的一个 $W^{k,p}$ 邻域通过
$\xi\mapsto\exp_u(\xi)$ 微分同胚于 $\mathcal M^*(A,\Sigma;J)$ 中 $u$ 的一个
$W^{k,p}$ 邻域。由命题~\ref{prop:3-1-1} 和
$\mathcal J_{\mathrm{reg}}(A,\Sigma)$ 的定义，微分
$d\mathcal F_u(0)=D_u$ 是满射。因为 $\mathcal F_u$ 是巴拿赫空间之间的光滑映射，
定理 A.3.3 表明，$\mathcal F_u^{-1}(0)$ 与零点的一个充分小邻域相交所得的集合，是维数为
\[
  n(2-2g)+2c_1(A)
\]
的光滑子流形。它在映射 $\xi\mapsto\exp_u(\xi)$ 下的像是
$W^{k,p}(\Sigma,M)$ 的光滑子流形，并与 $\mathcal M^*(A,\Sigma;J)$ 中 $u$ 的一个
邻域一致。因此 $\mathcal M^*(A,\Sigma;J)$ 是 $W^{k,p}(\Sigma,M)$ 的光滑子流形。
需要强调的是，如此得到的 $\mathcal M^*(A,\Sigma;J)$ 的坐标卡与 $k,p$ 无关，因为
$\mathcal F_u^{-1}(0)$ 中靠近零点的所有元素必定光滑（命题~\ref{prop:3-1-10}）。

为完成定理~\ref{thm:3-1-6} 第 \textup{(i)} 部分的证明，还须说明\BookRecordIndex{T0246}模空间具有典范定向。
先注意，
切空间 $T_u\mathcal M^*(A,\Sigma;J)$ 是 $D_u$ 的核。由式~\eqref{eq:3-1-7}，
$D_u$ 是算子
\[
  \xi\longmapsto(\widehat\nabla\xi)^{0,1}
\]
与形如 $\xi\mapsto T(\xi,\partial_J(u))$ 的算子之和。前者是一阶算子且与 $J$
交换，后者是零阶算子且与 $J$ 反交换（见注~\ref{rem:3-1-3}）。所以 $D_u$ 的核一般并不在 $J$
下不变，$J$ 也就未必在 $T_u\mathcal M^*(A,\Sigma;J)$ 上给出\BookRecordIndex{T0088}殆复结构。但是，把
$D_u$ 的第二部分乘以一个趋于零的常数，就可经一族\BookRecordIndex{T0128}弗雷德霍姆算子 $D_u^t$ 把 $D_u$
同伦到一个与 $J$ 交换的\BookRecordIndex{T0128}弗雷德霍姆算子，即复线性\BookRecordIndex{T0210}柯西--黎曼算子
\[
  \xi\longmapsto(\widehat\nabla\xi)^{0,1}.
\]

当复巴拿赫空间之间的\BookRecordIndex{T0128}弗雷德霍姆算子 $D:X\to Y$ 是复线性的时，它的行列式
\[
  \det(D)=\Lambda^{\max}(\ker D)\otimes
  \Lambda^{\max}(\ker D^*)
\]
带有自然定向（见第 A.2 节）。在这里，$D_u$ 的复反线性部分是紧算子，故\BookRecordIndex{T0162}行列式线
\[
  \det(D_u)=\Lambda^{\max}(\ker D_u)
\]
从 $D_u$ 的复线性部分继承一个自然定向。对所有
$u\in\mathcal M^*(A,\Sigma;J)$ 得到的这些 $\det(D_u)$ 的定向，确定了
$\mathcal M^*(A,\Sigma;J)$ 的定向。这就证明了定理~\ref{thm:3-1-6}\textup{(i)}。
\end{proof}

\begin{remark}\label{rem:3-2-4}
Donaldson 在 \srcite{81} 中曾用类似论证为 Yang--Mills \BookRecordIndex{T0246}模空间定向，Ruan
\srcite{342} 在当前语境中也使用了类似论证。McDuff 在 \srcite{254} 中采用了稍有不同的
论证方法，证明\BookRecordIndex{T0246}模空间 $\mathcal M^*(A,\Sigma;J)$ 的紧子集上存在典范稳定\BookRecordIndex{T0088}殆复结构\footnote{如果存在
某个整数 $k$，使切丛 $TX$ 与平凡 $k$ 维实向量丛的 Whitney 和
$TX\oplus\mathbb R^k$ 带有\BookRecordIndex{T0088}殆复结构，就称流形 $X$ 带有稳定\BookRecordIndex{T0088}殆复结构。}。
\end{remark}

\begin{remark}\label{rem:3-2-5}
若指标为零，则 $\mathcal M^*(A,\Sigma;J)$ 是离散集，每一点的定向都是一个符号
$\nu(u)=\pm1$。如命题 A.2.4 所述，这个符号可用穿越数解释。更明确地，选取一族光滑的
零阶（紧）算子
\[
  P^t:\Omega^0(\Sigma,u^*TM)
  \longrightarrow\Omega^{0,1}(\Sigma,u^*TM),
\]
使得 $P^0=0$，算子 $D_u^1:=D_u+P^1$ 复线性且双射，并且算子族
$t\mapsto D_u^t:=D_u+P^t$ 只有简单穿越。每个穿越都是使 $D_u^t$ 有非零（一维）核的
$t$ 值，而且
\[
  \nu(u)=(-1)^m,
\]
其中 $m$ 是穿越次数。因此，定向由算子族 $t\mapsto D_u^t$ 的模 $2$ \BookRecordIndex{T0273}谱流确定。注意，
这个符号的定义只适用于零维情形下 $\mathcal M^*(A,\Sigma;J)$ 中的单纯
\BookRecordIndex{T0003}$J$-全纯曲线。在某些情形（例如计数辛四维流形中的孤立嵌入环面）还必须计及\BookRecordIndex{T0113}多重覆盖
曲线；这时，确定它们对不变量的贡献要微妙得多（见 Taubes \srcite{389} 和
Ionel--Parker \srcite{198}）。
\end{remark}

\begin{remark}\label{rem:3-2-6}
若 $J$ 可积，则 $D_u$ 与 $J$ 交换，因而 $J$ 在
$\mathcal M^*(A,\Sigma;J)$ 上诱导一个（可积的）\BookRecordIndex{T0088}殆复结构。按定义，这一复结构与上面
定理~\ref{thm:3-1-6}\textup{(i)} 证明中描述的定向\BookRecordIndex{T0348}相容。特别地，若指标为零，且 $J$
正则、可积，则 $\mathcal M^*(A,\Sigma;J)$ 的每个元素对 $A$ 类 \BookRecordIndex{T0003}$J$-全纯曲线的计数都
贡献 $+1$。因此，在这个数为负的情形，不可能找到正则可积的 $J$（见例~\ref{ex:3-3-7}）。
\end{remark}

\begin{proof}[定理~\ref{thm:3-1-6}\textup{(ii)} 的证明]
现在须证明光滑正则\BookRecordIndex{T0088}殆复结构 $J$ 所成的集合
$\mathcal J_{\mathrm{reg}}(A,\Sigma)$ 在 $\mathcal J$ 中是\BookRecordIndex{T0300}剩余集。论证以投影
\[
  \pi:\mathcal M^*(A,\Sigma;\mathcal J^\ell)
  \longrightarrow\mathcal J^\ell
\]
的性质为基础。由命题~\ref{prop:3-2-1}（取 $k=1$），这是\BookRecordIndex{T0213}可分
$C^{\ell-1}$ \BookRecordIndex{T0053}巴拿赫流形之间的 $C^{\ell-1}$ 映射。切空间
\begin{multline*}
  T_{(u,J)}\mathcal M^*(A,\Sigma;\mathcal J^\ell)\\
  \subset W^{1,p}(\Sigma,u^*TM)
  \times C^\ell\bigl(M,\End(TM,J,\omega)\bigr)
\end{multline*}
由所有满足
\[
  D_u\xi+\frac12Y(u)\circ du\circ j_\Sigma=0
\]
的偶对 $(\xi,Y)$ 组成。此外，微分
\[
  d\pi(u,J):T_{(u,J)}\mathcal M^*(A,\Sigma;\mathcal J^\ell)
  \longrightarrow T_J\mathcal J^\ell
\]
就是投影 $(\xi,Y)\mapsto Y$。所以 $d\pi(u,J)$ 的核同构于 $D_u$ 的核；由
引理 A.3.6，它的余核也同构于 $D_u$ 的余核。因此 $d\pi(u,J)$ 是\BookRecordIndex{T0128}弗雷德霍姆算子，
并与 $D_u$ 有相同指标；而且 $d\pi(u,J)$ 是满射当且仅当 $D_u$ 是满射。

因此，$\pi$ 的\BookRecordIndex{T0409}正则值 $J$ 恰好是具有如下性质的\BookRecordIndex{T0088}殆复结构：对每条单纯 \BookRecordIndex{T0003}$J$-全纯曲线
$u\in\mathcal M^*(A,\Sigma;J)=\pi^{-1}(J)$，$D_u$ 都是满射。换言之，正则\BookRecordIndex{T0088}殆复
结构集合
\[
  \mathcal J_{\mathrm{reg}}^\ell(A,\Sigma)
  :=\bigl\{J\in\mathcal J^\ell\bigm|
  D_u\text{ 对所有 }u\in\mathcal M^*(A,\Sigma;J)
  \text{ 都是满射}\bigr\}
\]
正是 $\pi$ 的\BookRecordIndex{T0409}正则值集。由 萨德--斯梅尔定理（定理 A.5.1），这个集合在贝尔意义下
是\BookRecordIndex{T0300}剩余集（即可数个开稠密集的交）。这里用到了流形
$\mathcal M^*(A,\Sigma;\mathcal J^\ell)$ 与投影 $\pi$ 都属于 $C^{\ell-1}$；因此，只要
\[
  \ell-2\ge\operatorname{index}\pi
  =\operatorname{index}D_u
  =n(2-2g)+2c_1(A),
\]
就能应用 萨德--斯梅尔定理。于是，当 $\ell$ 充分大时，已证明
$\mathcal J_{\mathrm{reg}}^\ell(A,\Sigma)$ 在 $\mathcal J^\ell$ 的 $C^\ell$ 拓扑下
稠密。现在用 Taubes 的论证\BookRecordIndex{T0326}推出：$\mathcal J_{\mathrm{reg}}(A,\Sigma)$ 在
$\mathcal J$ 的 $C^\infty$ 拓扑下是\BookRecordIndex{T0300}剩余集。

考虑由所有如下光滑\BookRecordIndex{T0088}殆复结构 $J\in\mathcal J$ 组成的集合
\[
  \mathcal J_{\mathrm{reg},K}(A,\Sigma)\subset\mathcal J:
\]
对每条代表 $A$、满足
\begin{equation}\label{eq:3-2-3}
  \|du\|_{L^\infty}\le K
\end{equation}
并且存在 $z\in\Sigma$ 使
\begin{equation}\label{eq:3-2-4}
  \inf_{\zeta\ne z}
  \frac{d\bigl(u(z),u(\zeta)\bigr)}{d(z,\zeta)}
  \ge\frac1K
\end{equation}
的 \BookRecordIndex{T0003}$J$-全纯曲线 $u:\Sigma\to M$，算子 $D_u$ 都是满射。条件~\eqref{eq:3-2-4}
保证 $u$ 单纯；反过来，每条单纯 \BookRecordIndex{T0003}$J$-全纯曲线 $u:\Sigma\to M$ 都对某个 $K>0$
同时满足这两个条件。因此
\[
  \mathcal J_{\mathrm{reg}}(A,\Sigma)
  =\bigcap_{K>0}\mathcal J_{\mathrm{reg},K}(A,\Sigma).
\]
我们证明每个 $\mathcal J_{\mathrm{reg},K}(A,\Sigma)$ 都在 $\mathcal J$ 的
$C^\infty$ 拓扑下开且稠密。

先证 $\mathcal J_{\mathrm{reg},K}(A,\Sigma)$ 是开的，等价地，证明其补集闭。设
$J_\nu\notin\mathcal J_{\mathrm{reg},K}(A,\Sigma)$ 在 $C^\infty$ 拓扑下收敛于
$J\in\mathcal J$。对每个 $\nu$，存在一点 $z_\nu\in\Sigma$ 及一条满足
式~\eqref{eq:3-2-3}、式~\eqref{eq:3-2-4}（其中以 $z_\nu$ 代替 $z$）的
$J_\nu$-全纯曲线
\[
  u_\nu\in\mathcal M^*(A,\Sigma;J_\nu),
\]
并且 $D_{u_\nu}$ 不是满射。标准的\BookRecordIndex{T0330}椭圆自举论证表明，存在子列 $u_{\nu_i}$，连同所有
阶导数一起一致收敛于光滑 \BookRecordIndex{T0003}$J$-全纯曲线 $u:\Sigma\to M$（见定理 B.4.2）。再选取子列，
使 $z_{\nu_i}$ 收敛于 $z$。极限曲线 $u$ 对该点 $z$ 满足条件~\eqref{eq:3-2-3} 和
条件~\eqref{eq:3-2-4}；而且，由于各 $D_{u_{\nu_i}}$ 都不是满射，$D_u$ 也不可能是
满射。这说明 $J\notin\mathcal J_{\mathrm{reg},K}(A,\Sigma)$，故其补集在
$C^\infty$ 拓扑下闭。

再证 $\mathcal J_{\mathrm{reg},K}(A,\Sigma)$ 在 $\mathcal J$ 的 $C^\infty$ 拓扑下
稠密。首先注意
\[
  \mathcal J_{\mathrm{reg},K}(A,\Sigma)
  =\mathcal J_{\mathrm{reg},K}^\ell(A,\Sigma)\cap\mathcal J,
\]
其中 $\mathcal J_{\mathrm{reg},K}^\ell(A,\Sigma)\subset\mathcal J^\ell$ 由所有如下
$J\in\mathcal J^\ell(M,\omega)$ 组成：对每条代表 $A$、属于 $C^\ell$，并且对某个
$z\in\Sigma$ 满足式~\eqref{eq:3-2-3} 和式~\eqref{eq:3-2-4} 的 \BookRecordIndex{T0003}$J$-全纯曲线
$u:\Sigma\to M$，$D_u$ 都是满射。与上面相同的论证表明，
$\mathcal J_{\mathrm{reg},K}^\ell(A,\Sigma)$ 在 $\mathcal J^\ell$ 的 $C^\ell$ 拓扑下开。

任取 $J\in\mathcal J$。因为当 $\ell$ 很大时，
$\mathcal J_{\mathrm{reg}}^\ell(A,\Sigma)$ 在 $\mathcal J^\ell$ 中稠密，存在序列
\[
  J_\ell\in\mathcal J_{\mathrm{reg}}^\ell(A,\Sigma),
  \qquad \ell\ge\ell_0,
  \qquad
  \|J-J_\ell\|_{C^\ell}\le2^{-\ell}.
\]
由于 $J_\ell\in\mathcal J_{\mathrm{reg},K}^\ell(A,\Sigma)$，而后一个集合在
$C^\ell$ 拓扑下开，存在 $\varepsilon_\ell>0$，使得对每个 $J'\in\mathcal J^\ell$，
\[
  \|J'-J_\ell\|_{C^\ell}<\varepsilon_\ell
  \quad\Longrightarrow\quad
  J'\in\mathcal J_{\mathrm{reg},K}^\ell(A,\Sigma).
\]
选取任意光滑元素 $J_\ell'\in\mathcal J$，使得
\[
  \|J_\ell'-J_\ell\|_{C^\ell}
  <\min\{\varepsilon_\ell,2^{-\ell}\}.
\]
于是
\[
  J_\ell'\in
  \mathcal J_{\mathrm{reg},K}^\ell(A,\Sigma)\cap\mathcal J
  =\mathcal J_{\mathrm{reg},K}(A,\Sigma),
\]
而序列 $J_\ell'$ 在 $C^\infty$ 拓扑下收敛于 $J$。这就证明了所称的稠密性。因此，
$\mathcal J_{\mathrm{reg}}(A,\Sigma)$ 是可数个开稠密集
$\mathcal J_{\mathrm{reg},K}(A,\Sigma)$（$K\in\mathbb N$）的交，因而是\BookRecordIndex{T0300}剩余集。
定理~\ref{thm:3-1-6}\textup{(ii)} 得证。
\end{proof}

\begin{proof}[定理~\ref{thm:3-1-8} 的证明]
证明几乎逐字重复定理~\ref{thm:3-1-6} 的证明。记
$\mathcal J^\ell(J_0,J_1)$ 为从 $J_0$ 到 $J_1$ 的 $C^\ell$ 同伦
\[
  [0,1]\longrightarrow\mathcal J^\ell(M,\omega),
  \qquad \lambda\longmapsto J_\lambda
\]
组成的空间。考察\BookRecordIndex{T0268}普适模空间
\[
  \mathcal W^*(A,\Sigma;\mathcal J^\ell)
  \subset[0,1]\times\mathcal B^{1,p}
  \times\mathcal J^\ell(J_0,J_1),
\]
它由所有三元组 $(\lambda,u,J)$ 组成，其中 $\lambda\in[0,1]$、
$J\in\mathcal J^\ell(J_0,J_1)$，且
$(\lambda,u)\in\mathcal W^*(A,\Sigma;J)$。与命题~\ref{prop:3-2-1} 和
定理~\ref{thm:3-1-6} 的证明相同，可以证明
$\mathcal W^*(A,\Sigma;\mathcal J^\ell)$ 是
$[0,1]\times\mathcal B^{1,p}\times\mathcal J^\ell(J_0,J_1)$ 的
$C^{\ell-1}$ 巴拿赫子流形。它的边界是它与
$\{0,1\}\times\mathcal B^{1,p}\times\mathcal J^\ell(J_0,J_1)$ 的交。投影
\begin{equation}\label{eq:3-2-5}
  \mathcal W^*(A,\Sigma;\mathcal J^\ell)
  \longrightarrow\mathcal J^\ell(J_0,J_1),
  \qquad
  (\lambda,u,J)\longmapsto J
\end{equation}
仍是 $C^{\ell-1}$ 弗雷德霍姆映射，其\BookRecordIndex{T0409}正则值组成 $C^\ell$ 类\BookRecordIndex{T0407}正则同伦的集合
$\mathcal J_{\mathrm{reg}}^\ell(A,\Sigma;J_0,J_1)$。由 萨德--斯梅尔定理，\BookRecordIndex{T0409}正则值集
在 $\mathcal J^\ell(J_0,J_1)$ 中稠密。从光滑情形约化到 $C^\ell$ 情形的方法与
定理~\ref{thm:3-1-6} 的证明相同。因此，
$\mathcal J_{\mathrm{reg}}(A,\Sigma;J_0,J_1)$ 是
$\mathcal J(J_0,J_1)$ 的剩余子集。

若光滑\BookRecordIndex{T0407}正则同伦
$\{J_\lambda\}_\lambda\in\mathcal J_{\mathrm{reg}}(A,\Sigma;J_0,J_1)$，则由\BookRecordIndex{T0329}椭圆
正则性（命题~\ref{prop:3-1-10}），光滑\BookRecordIndex{T0246}模空间
$\mathcal W^*(A,\Sigma;\{J_\lambda\}_\lambda)$ 对每个 $\ell$ 都是投影
\eqref{eq:3-2-5} 在 $\{J_\lambda\}_\lambda$ 上的原像。因此，对每个整数 $\ell$，它
都是 $[0,1]\times\mathcal B^{1,p}$ 的 $C^{\ell-1}$ 子流形，从而是光滑子流形；其边界
是它与 $\{0,1\}\times\mathcal B^{1,p}$ 的交。这就证明了
定理~\ref{thm:3-1-8}。
\end{proof}

\begin{remark}\label{rem:3-2-7}
定理~\ref{thm:3-1-6}\textup{(ii)} 的证明要点，是对每个 $K$ 证明
$\mathcal J_{\mathrm{reg},K}(A,\Sigma)$ 在 $\mathcal J$ 中稠密。上面的证明用
$C^\ell$ \BookRecordIndex{T0053}巴拿赫流形上的 萨德--斯梅尔定理来做到这一点。也可采用 Floer 的论证，始终
留在 $C^\infty$ 范畴中。更确切地，可以考虑由所有满足
\[
  \|Y\|_\varepsilon
  :=\sum_{\ell=1}^\infty
  \varepsilon_\ell\|Y\|_{C^\ell}<\infty
\]
的截面 $Y\in C^\infty(M,\End(TM,J,\omega))$ 组成的一个适当巴拿赫空间
\[
  C_\varepsilon^\infty\bigl(M,\End(TM,J,\omega)\bigr)
  \subset C^\infty\bigl(M,\End(TM,J,\omega)\bigr).
\]
容易看出，如果序列 $\varepsilon_\ell$ 足够快地趋于零，那么在
$C_\varepsilon^\infty(M,\End(TM,J,\omega))$ 中，存在支集含于任意点
$x_0\in M$ 的任意小邻域、且在 $x_0$ 处可取任意值的截面 $Y$。（细节见 Floer
\srcite{114}。）因此，在命题~\ref{prop:3-2-1} 的证明中，可以把 $\mathcal J^\ell$
换成所有形如
\[
  J=J_0\exp(-J_0Y),
  \qquad
  Y\in C_\varepsilon^\infty
  \bigl(M,\End(TM,J_0,\omega)\bigr)
\]
的\BookRecordIndex{T0088}殆复结构组成的空间 $\mathcal J_\varepsilon$。由于 $\mathcal J_\varepsilon$ 是光滑\BookRecordIndex{T0213}可分
\BookRecordIndex{T0053}巴拿赫流形，萨德--斯梅尔定理适用。于是，与定理~\ref{thm:3-1-6}\textup{(ii)} 的
证明相同，每个 $J_0\in\mathcal J$ 都能在 $C^\infty$ 拓扑下由一列正则\BookRecordIndex{T0088}殆复结构逼近。
\end{remark}

\begin{remark}\label{rem:3-2-8}
定理~\ref{thm:3-1-6} 的证明表明，光滑\BookRecordIndex{T0268}普适模空间
$\mathcal M^*(A,\Sigma;\mathcal J)$ 中一点 $(u,J)$ 按定义~\ref{def:3-1-5} 是正则的，
当且仅当投影
\[
  d\pi(u,J):T_{(u,J)}\mathcal M^*(A,\Sigma;\mathcal J)
  \longrightarrow T_J\mathcal J
\]
是满射。由\BookRecordIndex{T0380}隐函数定理，任意从 $J_0=J$ 出发的光滑道路
\[
  [0,1]\longrightarrow\mathcal J,
  \qquad t\longmapsto J_t
\]
都可在某个区间 $[0,\varepsilon)$ 上提升为\BookRecordIndex{T0268}普适模空间中从 $u_0=u$ 出发的道路
\[
  [0,\varepsilon)\longrightarrow
  \mathcal M^*(A,\Sigma;\mathcal J),
  \qquad t\longmapsto(u_t,J_t).
\]
因此，任意正则曲线 $(u,J)$ 在扰动 $J$ 时都会持续存在。反之，若 $(u,J)$ 不正则，
扰动 $J$ 后可能不存在邻近曲线（见例~\ref{ex:3-3-6}）。
\end{remark}
